\documentclass[../main.tex]{subfiles}
\graphicspath{{\subfix{../images/}}}
\begin{document}


\section{Konečně-rozměrné dynamické systémy generované ODR}


\subsection{Úloha 1}\label{sec:uloha1}

\begin{equation}
    \frac{d x}{d t} = f(x) , x(0) = x_ini 
\end{equation}
,kde 
$\Gamma \subset \mathbb{R}^N$ je oblast, $f: \Gamma \mapsto \mathbb{R}^n, f\in C^{(3)}(\Gamma)$ a $x_{ini} \in \Gamma$

\todo{Stuff}




\begin{theorem}[Liouville]
    Nechť platí pro \ref{sec:uloha1} předpoklady (z)\todo{change to link} a navíc 
    $f\in C^{(3)} (\Gamma)$ a $M\subset\Gamma$ je měřitelná a $\mu(M) < \infty$.

    Pak\begin{equation}
        \frac{d}{dt} \left(\mu (\Phi_t (M))\right)|_{t=0} = \int_M \div f(x) dx
    \end{equation}
\end{theorem}

\begin{proof}
    \begin{equation}
        \mu(\Phi_t(M)) = \int_{\Phi_t(M)} 1 dy = \{\text{substituce }y =\Phi_t(x) \} = \int_M 1 |\mathcal{Y}|dx
    \end{equation}

    Výpočet jakobiánu

    \begin{equation}
        |\mathcal{Y}| = 
            \begin{vmatrix}
            \frac{\partial (\Phi_t(x))^1}{\partial x^1} & \dots & \frac{\partial (\Phi_t(x))^n}{\partial x^1} \\
            \vdots                                      &       & \vdots                                      \\
            \frac{\partial (\Phi_t(x))^1}{\partial x^n} & \dots & \frac{\partial (\Phi_t(x))^n}{\partial x^n}
            \end{vmatrix}
    \end{equation}

    Pro $t$ z okolí $0$ a $\hat{x} \in \Gamma$:
    \begin{equation}\label{eq:zelenanula}
        \Phi_t(\hat{x}) = \Phi_0(\hat{x}) + t \frac{\partial}{\partial t } (\Phi_t(\hat{x}))|_{t=0} + \frac{1}{2}t^2 \frac{\partial^2}{\partial t^2} (\Phi_t (\hat{x})) |_{t=0} + R_3(t, \Phi_t(\hat{x}))
    \end{equation} 

    zde
    \begin{equation}
        \frac{\partial}{\partial t }(\Phi_t(\hat{x })) = f(\Phi_t (\hat{x}))
    \end{equation}
    a 
    \begin{equation}
        \frac{\partial^2}{\partial t^2} (\Phi_t (\hat{x})) = \frac{\partial}{\partial t}(f(\Phi_t (\hat{x}))) = f'(\Phi_t (\hat{x})) \frac{\partial}{\partial t} (\Phi_t (\hat{x})) = f'(\Phi_t (\hat{x})) f(\Phi_t (\hat{x}))
    \end{equation}

    Rovnici \eqref{eq:zelenanula} derivijeme podle složek $\hat{x}$

\end{proof}










\end{document}